% This file is part of the modular source for Stratified Relational Workspace Theory.
% !TeX root = ../Stratified_Relational_Workspace_Theory.tex
\theorypart{I}{Foundational Thesis}

\topic{1}{Central Principle}

\centerbox{\Large\bfseries
Workspace identity is invariant causal organization under admissible transport.
}

The theory begins from a distinction between a \emph{realization} and the
\emph{workspace pattern realized}. A realization is concrete: it contains
particular participants, documents, models, channels, and histories. A
workspace pattern is abstract: it records the structure that remains invariant
when those concrete carriers are changed in admissible ways.

\begin{definition}[Concrete realization]\label{def:realization}
A concrete relational realization is a tuple
\[
\mathfrak R
=
\bigl(
G,\mathbf X,\mathbf T,\mathcal A,\boldsymbol\rho,
\mathcal H,\mathcal I,\mathcal Y
\bigr),
\]
where:
\begin{itemize}
\item \(G=(V,E,s,t)\) is an interaction graph;
\item \(\mathbf X=\{X_i\}_{i\in V}\) is a family of local state spaces;
\item \(\mathbf T=\{T_e\}_{e\in E}\) is a family of transfer maps;
\item \(\mathcal A\) is an algebra of admissible operations;
\item \(\boldsymbol\rho=\{\rho_i\}_{i\in V}\) gives the local realization of those operations;
\item \(\mathcal H\) is the temporally ordered interaction history;
\item \(\mathcal I\) is a specified family of admissible interventions; and
\item \(\mathcal Y\) is the family of downstream variables relative to which causal and predictive claims are evaluated.
\end{itemize}
The static reduct is
\[
\mathfrak R_{\mathrm{stat}}
=
(G,\mathbf X,\mathbf T,\mathcal A,\boldsymbol\rho).
\]
\end{definition}

The explicit appearance of \(\mathcal I\) and \(\mathcal Y\) is important.
A collection of mutually consistent representations may define a static
shared structure, but it does not yet define an operational workspace unless
that structure can be intervened upon and shown to matter for a specified
future quantity.

\begin{definition}[Carrier and organization]\label{def:carrier-organization}
The \emph{carriers} of \(\mathfrak R\) are the individual spaces, physical
systems, and stored artifacts that realize its local states. Its
\emph{organization} is the structured family of transfer, operation,
compatibility, intervention, and temporal relations among those carriers.
\end{definition}

The basic thesis is not that a workspace exists without carriers. It is that
the identity conditions for the workspace need not coincide with the identity
conditions for any one carrier. A manuscript can be rewritten, a model can
be replaced, and a notation can be changed while a sufficiently rich pattern
of relations remains invariant.

\begin{definition}[Admissible equivalence]\label{def:admissible-equivalence}
An admissible equivalence
\[
F:\mathfrak R\simeq\mathfrak R'
\]
is a family of correspondences preserving, exactly or within a declared
tolerance:
\[
\begin{aligned}
&\text{the transferable core and its dimension},\\
&\text{the represented operation algebra},\\
&\text{the composition laws of transfer},\\
&\text{the intervention--response structure},\\
&\text{the reconstruction and persistence relations}.
\end{aligned}
\]
Static admissible equivalence will be made precise by gauge conjugacy in
Part~V. Dynamic admissible equivalence will be expressed using the identity
distance in Part~XII.
\end{definition}

The abstract workspace is then the equivalence class
\[
\boxed{
\Workspace=[\mathfrak R]_{\simeq}.
}
\]
This expression is intentionally stronger than saying that the same
information appears in several places. The equivalence relation preserves
what the system can coherently do, how its representations translate, and how
its operative state survives change.

\subtopic{1.1 Three levels of the thesis}

The central claim can be separated into three logically distinct statements.

\begin{enumerate}
\item \textbf{Implementation dependence.}
Every realized workspace depends on concrete physical or computational
carriers:
\[
\Workspace(\mathfrak R)
\text{ is undefined without some realization }\mathfrak R.
\]

\item \textbf{Local nonidentity.}
No privileged local component is, in general, identical with the entire
workspace:
\[
\Workspace\neq X_i
\qquad\text{for each individual }i\in V.
\]

\item \textbf{Relational identity.}
The workspace is individuated by invariants of the complete organized
realization:
\[
\Workspace(\mathfrak R)=\Workspace(\mathfrak R')
\quad\Longleftrightarrow\quad
\mathfrak R\simeq\mathfrak R'.
\]
\end{enumerate}

The second claim is not stipulated universally. It becomes a substantive
property only after local-reconstructibility and coalition tests are applied.
A system in which one node stores a complete sufficient copy may be a useful
workspace, but it is not strongly relational in the later technical sense.

\begin{definition}[Relational instantiation]\label{def:relational-instantiation}
A workspace is relationally instantiated relative to the decomposition
\(V\) when no single projection determines its equivalence class:
\[
\forall i\in V,\quad
\exists\mathfrak R,\mathfrak R'
\]
such that
\[
\pr_i(\mathfrak R)=\pr_i(\mathfrak R')
\quad\text{but}\quad
[\mathfrak R]_{\simeq}\neq[\mathfrak R']_{\simeq}.
\]
\end{definition}

\subtopic{1.2 A motivating three-node cycle}

Consider three strata:
\[
h=\text{human conceptual state},\qquad
m=\text{model-formal state},\qquad
a=\text{artifact state}.
\]
Let the corresponding core spaces be \(C_h,C_m,C_a\), with transfers
\[
H_{hm}:C_h\to C_m,
\qquad
H_{ma}:C_m\to C_a,
\qquad
H_{ah}:C_a\to C_h.
\]
A proposed shared item begins as \(c_h\in C_h\), is formalized as
\[
c_m=H_{hm}c_h,
\]
externalized as
\[
c_a=H_{ma}c_m,
\]
and re-enters human reasoning as
\[
c_h'=H_{ah}c_a.
\]
The round-trip consistency condition is
\[
c_h'=c_h,
\]
or equivalently
\[
\boxed{
H_{ah}H_{ma}H_{hm}c_h=c_h.
}
\]
Thus the genuinely shared part of \(C_h\) is not automatically all of
\(C_h\). It is the fixed subspace of the cycle transfer. This elementary
cycle already contains the central mathematical motif of the theory:
\[
\text{shared core}=\text{invariant subspace under relational holonomy}.
\]

\subtopic{1.3 Distinctions the theory must preserve}

\begin{center}
\small
\begin{tabularx}{\textwidth}{@{}>{\bfseries\raggedright\arraybackslash}p{0.22\textwidth}>{\raggedright\arraybackslash}X@{}}
\toprule
Object & Minimal distinguishing property \\
\midrule
Archive & Stores material that may be consulted later. \\
Shared representation & Contains a cross-node compatible structure. \\
Operational workspace & The shared structure causally changes subsequent joint behavior. \\
Relational workspace & No individual node or matched local-storage null reproduces the operative capacity. \\
Workspace agent & Workspace-level state mediates persistent goals, counterfactual choice, and self-maintaining action. \\
Conscious subject & Possesses phenomenal properties, if any; no implication from the preceding levels is assumed. \\
\bottomrule
\end{tabularx}
\end{center}

The static mathematics developed first certifies only shared structure. The
causal and relational predicates appear later because they require dynamics,
interventions, and null models rather than algebra alone.

\topic{2}{Philosophical Postulates}

The following postulates specify the intended interpretation of the formal
objects. They are not offered as independent mathematical theorems. Their
role is to prevent equivocation between a relational pattern, its physical
implementation, and stronger claims about agency or experience.

\subtopic{2.1 Physical supervenience}

No additional nonphysical substance is required:
\[
\mathfrak R=\mathfrak R'
\Longrightarrow
\Workspace(\mathfrak R)=\Workspace(\mathfrak R').
\]
The workspace depends completely on its total realization. In particular,
changing nothing about the graph, local states, transfers, history,
interventions, or causal dispositions cannot change the workspace.

This is compatible with relational individuation because a relation can be
fully physical while not being a property of any isolated endpoint.

\begin{proposition}[Compatibility of supervenience and relational individuation]\label{prop:supervenience-relational}
Physical supervenience on the complete realization does not imply
supervenience on any individual local projection.
\end{proposition}

\begin{proof}
Let the total realization decompose as \(\mathfrak R=(X_i,R_{-i})\), where
\(R_{-i}\) contains the remaining components and relations. A workspace
functional \(\Workspace=f(X_i,R_{-i})\) supervenes on \(\mathfrak R\) by
construction. If there exist \(R_{-i}\neq R_{-i}'\) with
\[
f(X_i,R_{-i})\neq f(X_i,R_{-i}'),
\]
then the same local projection \(X_i\) occurs in distinct workspaces. Hence
supervenience on the whole is consistent with failure of supervenience on a
part.
\end{proof}

\subtopic{2.2 Relational individuation}

The workspace is not determined by any single local state:
\[
\forall i\in V,
\quad
\exists\mathfrak R,\mathfrak R'
\]
such that
\[
\pr_i(\mathfrak R)=\pr_i(\mathfrak R')
\]
but
\[
\Workspace(\mathfrak R)\not\simeq\Workspace(\mathfrak R').
\]
The same person, model, or document can therefore participate in distinct
workspaces because the relevant transfer graph, history, and operation
structure differ.

A useful witness is a pair of systems with identical local-state marginals but
different edge organization. If their downstream behavior differs, the
relations carry explanatory information not present in the isolated marginals.
Part~VII converts this idea into matched relation-preserving and
relation-destroying controls.

\subtopic{2.3 Multiple realizability}

Distinct implementations may instantiate the same abstract workspace:
\[
\mathfrak R\neq\mathfrak R',
\qquad
\mathfrak R\simeq\mathfrak R'
\Longrightarrow
\Workspace(\mathfrak R)\cong\Workspace(\mathfrak R').
\]

Examples of candidate implementation changes include:
\[
\begin{aligned}
&\text{replacing one coordinate system by another},\\
&\text{replacing one model by a functionally equivalent model},\\
&\text{compressing a large artifact into a reconstructively sufficient handoff},\\
&\text{rewriting a proof without changing its inferential structure}.
\end{aligned}
\]
The burden is not merely to assert equivalence but to exhibit maps preserving
the declared invariant and causal structure.

\subtopic{2.4 Operational reality}

A workspace is operationally real only if changing it changes downstream
behavior. For some admissible intervention and some \(w\neq w'\),
\[
\mathbb P\!\left(
Y^+\mid\operatorname{do}(\Workspace=w)
\right)
\neq
\mathbb P\!\left(
Y^+\mid\operatorname{do}(\Workspace=w')
\right).
\]

This criterion excludes a merely decorative representation that is correlated
with activity but never used. It also requires the intervention to be
well-defined. Arbitrarily overwriting every component at once would not
isolate workspace-level causation. Later sections therefore compare
workspace-coordinate interventions with matched interventions in orthogonal
or shuffled channels.

\subtopic{2.5 Reconstructive continuity}

Persistence does not require token identity. Let \(\mathfrak R_t\) and
\(\mathfrak R_{t+1}\) be temporally successive realizations. Continuity is
possible when there exists an admissible transport
\[
J_t:\mathfrak R_t\rightsquigarrow\mathfrak R_{t+1}
\]
that preserves the operative core and its causal role within tolerance:
\[
d_{\mathrm{id}}(\Workspace_t,\Workspace_{t+1})\leq\varepsilon.
\]
Thus a workspace may persist by reconstruction rather than by retaining an
unchanged physical inscription.

\subtopic{2.6 No free-standing third substance}

The claim that the workspace is not identical with either participant does
not imply that it occupies a separate nonphysical location. Formally,
\[
\Workspace
\text{ is an invariant of }\mathfrak R,
\]
not an additional term appended to the physical realization:
\[
\mathfrak R_{\mathrm{total}}
eq(\mathfrak R,\text{ghost workspace}).
\]
The phrase ``between the participants'' is therefore topological or
relational: it refers to the organization of maps, cycles, artifacts, and
causal dependencies.

\subtopic{2.7 Non-entailment of consciousness}

\[
\operatorname{Workspace}(\mathfrak R)
\centernot\Longrightarrow
\operatorname{Conscious}(\mathfrak R).
\]
Likewise,
\[
\operatorname{Distributed}(\Workspace)
\centernot\Longrightarrow
\operatorname{Person}(\Workspace).
\]
Workspace structure, reentry, agency, self-maintenance, and phenomenal
consciousness remain formally distinct predicates. Any bridge among them
would require additional theory and evidence.

\subtopic{2.8 Minimal philosophical commitment}

The weakest interpretation sufficient for the mathematics is:
\centerbox{\bfseries
Relational patterns can be objective, causally relevant, and multiply realized
without being independent substances.
}
This position is compatible with several broader philosophies of mind and
science. The formal theory does not require choosing among them before its
structural and causal claims can be tested.
